% 第17章17.4.1极大匹配与最大匹配：左侧 a,b，右侧 1,2，允许边 (a,1),(a,2),(b,1)。
% 左侧 b 在上、a 在下，使三条边不交叉。{(a,1)} 不能再加边（极大），但 {(a,2),(b,1)} 有两条边（最大，也是完美匹配）。
\begin{tikzpicture}[line width=.8pt,font=\fontsize{8}{10}\selectfont,
  v/.style={circle,draw=FigureStroke,fill=FigureBlueFill,line width=.8pt,minimum size=6.5mm,inner sep=0pt,text=FigureInk},
  e/.style={draw=FigureMuted,line width=.8pt},
  m/.style={draw=FigureBlue,line width=1.8pt}]
% Panel labels share one baseline, 3 mm below the full row of content.
\begin{scope}[local bounding box=panel-content]
  \foreach \dx/\k in {0/A,3.6/B}{
    \node[v] (a\k) at (\dx,-.6) {$a$};
    \node[v] (b\k) at (\dx,.6) {$b$};
    \node[v] (x\k) at (\dx+1.8,.6) {$1$};
    \node[v] (y\k) at (\dx+1.8,-.6) {$2$};
    \draw[e] (a\k)--(x\k); \draw[e] (a\k)--(y\k); \draw[e] (b\k)--(x\k);
  }
  \draw[m] (aA)--(xA);
  \draw[m] (aB)--(yB); \draw[m] (bB)--(xB);
  \coordinate[overlay] (panel-a-center) at (.9,-1.45);
  \coordinate[overlay] (panel-b-center) at (4.5,-1.45);
\end{scope}
\node[font=\figurefont\fontsize{8}{10}\selectfont,text=FigureInk,anchor=north,inner sep=0pt,outer sep=0pt] at ([yshift=-3mm]panel-a-center |- panel-content.south) {(a)};
\node[font=\figurefont\fontsize{8}{10}\selectfont,text=FigureInk,anchor=north,inner sep=0pt,outer sep=0pt] at ([yshift=-3mm]panel-b-center |- panel-content.south) {(b)};
\end{tikzpicture}
